% Options for packages loaded elsewhere
\PassOptionsToPackage{unicode}{hyperref}
\PassOptionsToPackage{hyphens}{url}
\documentclass[
  english,
  11pt,
  a4paper,
]{article}
\usepackage{xcolor}
\usepackage[margin=23mm]{geometry}
\usepackage{amsmath,amssymb}
\setcounter{secnumdepth}{-\maxdimen} % remove section numbering
\usepackage{iftex}
\ifPDFTeX
  \usepackage[T1]{fontenc}
  \usepackage[utf8]{inputenc}
  \usepackage{textcomp} % provide euro and other symbols
\else % if luatex or xetex
  \usepackage{unicode-math} % this also loads fontspec
  \defaultfontfeatures{Scale=MatchLowercase}
  \defaultfontfeatures[\rmfamily]{Ligatures=TeX,Scale=1}
\fi
\usepackage{lmodern}
\ifPDFTeX\else
  % xetex/luatex font selection
  \setmainfont[]{DejaVu Serif}
  \setsansfont[]{DejaVu Sans}
  \setmonofont[Scale=0.82]{DejaVu Sans Mono}
\fi
% Use upquote if available, for straight quotes in verbatim environments
\IfFileExists{upquote.sty}{\usepackage{upquote}}{}
\IfFileExists{microtype.sty}{% use microtype if available
  \usepackage[]{microtype}
  \UseMicrotypeSet[protrusion]{basicmath} % disable protrusion for tt fonts
}{}
\makeatletter
\@ifundefined{KOMAClassName}{% if non-KOMA class
  \IfFileExists{parskip.sty}{%
    \usepackage{parskip}
  }{% else
    \setlength{\parindent}{0pt}
    \setlength{\parskip}{6pt plus 2pt minus 1pt}}
}{% if KOMA class
  \KOMAoptions{parskip=half}}
\makeatother
\usepackage{color}
\usepackage{fancyvrb}
\newcommand{\VerbBar}{|}
\newcommand{\VERB}{\Verb[commandchars=\\\{\}]}
\DefineVerbatimEnvironment{Highlighting}{Verbatim}{commandchars=\\\{\}}
% Add ',fontsize=\small' for more characters per line
\newenvironment{Shaded}{}{}
\newcommand{\AlertTok}[1]{\textcolor[rgb]{1.00,0.00,0.00}{\textbf{#1}}}
\newcommand{\AnnotationTok}[1]{\textcolor[rgb]{0.38,0.63,0.69}{\textbf{\textit{#1}}}}
\newcommand{\AttributeTok}[1]{\textcolor[rgb]{0.49,0.56,0.16}{#1}}
\newcommand{\BaseNTok}[1]{\textcolor[rgb]{0.25,0.63,0.44}{#1}}
\newcommand{\BuiltInTok}[1]{\textcolor[rgb]{0.00,0.50,0.00}{#1}}
\newcommand{\CharTok}[1]{\textcolor[rgb]{0.25,0.44,0.63}{#1}}
\newcommand{\CommentTok}[1]{\textcolor[rgb]{0.38,0.63,0.69}{\textit{#1}}}
\newcommand{\CommentVarTok}[1]{\textcolor[rgb]{0.38,0.63,0.69}{\textbf{\textit{#1}}}}
\newcommand{\ConstantTok}[1]{\textcolor[rgb]{0.53,0.00,0.00}{#1}}
\newcommand{\ControlFlowTok}[1]{\textcolor[rgb]{0.00,0.44,0.13}{\textbf{#1}}}
\newcommand{\DataTypeTok}[1]{\textcolor[rgb]{0.56,0.13,0.00}{#1}}
\newcommand{\DecValTok}[1]{\textcolor[rgb]{0.25,0.63,0.44}{#1}}
\newcommand{\DocumentationTok}[1]{\textcolor[rgb]{0.73,0.13,0.13}{\textit{#1}}}
\newcommand{\ErrorTok}[1]{\textcolor[rgb]{1.00,0.00,0.00}{\textbf{#1}}}
\newcommand{\ExtensionTok}[1]{#1}
\newcommand{\FloatTok}[1]{\textcolor[rgb]{0.25,0.63,0.44}{#1}}
\newcommand{\FunctionTok}[1]{\textcolor[rgb]{0.02,0.16,0.49}{#1}}
\newcommand{\ImportTok}[1]{\textcolor[rgb]{0.00,0.50,0.00}{\textbf{#1}}}
\newcommand{\InformationTok}[1]{\textcolor[rgb]{0.38,0.63,0.69}{\textbf{\textit{#1}}}}
\newcommand{\KeywordTok}[1]{\textcolor[rgb]{0.00,0.44,0.13}{\textbf{#1}}}
\newcommand{\NormalTok}[1]{#1}
\newcommand{\OperatorTok}[1]{\textcolor[rgb]{0.40,0.40,0.40}{#1}}
\newcommand{\OtherTok}[1]{\textcolor[rgb]{0.00,0.44,0.13}{#1}}
\newcommand{\PreprocessorTok}[1]{\textcolor[rgb]{0.74,0.48,0.00}{#1}}
\newcommand{\RegionMarkerTok}[1]{#1}
\newcommand{\SpecialCharTok}[1]{\textcolor[rgb]{0.25,0.44,0.63}{#1}}
\newcommand{\SpecialStringTok}[1]{\textcolor[rgb]{0.73,0.40,0.53}{#1}}
\newcommand{\StringTok}[1]{\textcolor[rgb]{0.25,0.44,0.63}{#1}}
\newcommand{\VariableTok}[1]{\textcolor[rgb]{0.10,0.09,0.49}{#1}}
\newcommand{\VerbatimStringTok}[1]{\textcolor[rgb]{0.25,0.44,0.63}{#1}}
\newcommand{\WarningTok}[1]{\textcolor[rgb]{0.38,0.63,0.69}{\textbf{\textit{#1}}}}
\usepackage{longtable,booktabs,array}
\usepackage{caption}
\captionsetup[table]{skip=6pt}
\newcounter{none} % for unnumbered tables
\usepackage{calc} % for calculating minipage widths
% Correct order of tables after \paragraph or \subparagraph
\usepackage{etoolbox}
\makeatletter
\patchcmd\longtable{\par}{\if@noskipsec\mbox{}\fi\par}{}{}
\makeatother
% Allow footnotes in longtable head/foot
\IfFileExists{footnotehyper.sty}{\usepackage{footnotehyper}}{\usepackage{footnote}}
\makesavenoteenv{longtable}
\ifLuaTeX
\usepackage[bidi=basic,shorthands=off]{babel}
\else
\usepackage[bidi=default,shorthands=off]{babel}
\fi
\ifPDFTeX
\else
\babelfont{rm}[]{DejaVu Serif}
\fi
\ifLuaTeX
  \usepackage{selnolig} % disable illegal ligatures
\fi
\setlength{\emergencystretch}{3em} % prevent overfull lines
\providecommand{\tightlist}{%
  \setlength{\itemsep}{0pt}\setlength{\parskip}{0pt}}
\usepackage{fancyhdr}
\usepackage{titlesec}
\usepackage{needspace}
\definecolor{reviewblue}{HTML}{174E67}
\AtBeginDocument{\hypersetup{colorlinks=true,linkcolor=reviewblue,urlcolor=reviewblue,citecolor=reviewblue}}
\pagestyle{fancy}
\fancyhf{}
\fancyhead[L]{\small\sffamily CHAMP technical assessment}
\fancyhead[R]{\small\sffamily Version 1.0.1}
\fancyfoot[L]{\footnotesize\sffamily Mounir IDRASSI}
\fancyfoot[R]{\footnotesize\sffamily\thepage}
\setlength{\headheight}{15pt}
\renewcommand{\headrulewidth}{0.3pt}
\renewcommand{\footrulewidth}{0pt}
\fancypagestyle{plain}{\fancyhf{}\fancyfoot[L]{\footnotesize\sffamily Mounir IDRASSI}\fancyfoot[R]{\footnotesize\sffamily\thepage}\renewcommand{\headrulewidth}{0pt}}
\titleformat{\section}{\large\sffamily\bfseries\color{reviewblue}}{}{0pt}{}
\titlespacing*{\section}{0pt}{2.1ex plus 1ex minus .2ex}{1ex}
\AtBeginEnvironment{longtable}{\small}
\setlength{\emergencystretch}{3em}
\clubpenalty=10000
\widowpenalty=10000
\displaywidowpenalty=10000
\brokenpenalty=10000
\predisplaypenalty=10000
\usepackage{bookmark}
\IfFileExists{xurl.sty}{\usepackage{xurl}}{} % add URL line breaks if available
\urlstyle{same}
% fallback for those not using the hyperref driver hyperxmp:
\makeatletter
\@ifundefined{xmpquote}{\newcommand{\xmpquote}[1]{#1}}{}
\makeatother
\hypersetup{
  pdftitle={CHAMP: technical assessment of the submitted specification},
  pdfauthor={Mounir IDRASSI (mounir@amcrypto.jp)},
  pdflang={en},
  hidelinks,
  pdfcreator={LaTeX via pandoc}}

\title{CHAMP: technical assessment of the submitted specification}
\usepackage{etoolbox}
\makeatletter
\providecommand{\subtitle}[1]{% add subtitle to \maketitle
  \apptocmd{\@title}{\par {\large #1 \par}}{}{}
}
\makeatother
\subtitle{Independent cryptographic review · Version 1.0.1}
\author{Mounir IDRASSI
(\href{mailto:mounir@amcrypto.jp}{\nolinkurl{mounir@amcrypto.jp}})}
\date{22 September 2026}

\begin{document}
\maketitle

\textbf{Scope and version.} This report assesses the supplied ten-page
\emph{CHAMP Algorithm Specifications}, its four C implementations, test
vectors and statistical reports. All references to specification
sections and pages concern that document. Its SHA-256 is recorded in the
source inventory and in the references below.

A distinct twelve-page
\href{https://shpilrain.ccny.cuny.edu/Cayley\%20hash\%20paper.pdf}{author-hosted
preprint}, retrieved on 21 September 2026, cites Mullan--Tsaban in §6.2.
The supplied ten-page specification does not. The criticism here
concerns the absence of an analysis of the subgroup strategy for the
implemented determinant-2 pair, not a claim that the authors have never
cited that work. This report is not a comprehensive review of the
separate preprint; both document hashes are recorded in
\url{PROVENANCE.md}.

\textbf{Independent preparation.} This review's analysis and experiments
were completed independently of
\href{https://ngcc.dev/reports/hash-04.html}{Saarinen's report} and
comments signed
\href{https://list.niccs.org.cn/archives/list/crypthashforum@list.niccs.org.cn/message/QNMIAAKIPHG5FEK7QZEE2ZYBOWQHCZUE/}{Tsinghua
Hash Lab} and
\href{https://list.niccs.org.cn/archives/list/crypthashforum@list.niccs.org.cn/message/6LSOFTTMTDWFWWRAZDREDYJNG5NXDVEC/}{ISCAS}.
The author learned of these sources on 22 September 2026, after
completing the review and reproducibility package. Their citations
acknowledge overlap. The subgroup method is credited to Mullan--Tsaban
in §2.

\textbf{Additional contributions.} Relative to the cited sources, the
independently developed additions are the explicit two-stage subgroup
search with twelve verified reduced-parameter collision certificates
(§2); demonstrated prefix/suffix and three-query sandwich MAC forgeries,
and related-digest recovery (§4); a 40-bit preimage recovery using full
CHAMP-1024 parameters (§5); and digest-length, initialization and timing
findings (§7). The determinant bound, collision-search direction,
identity example, serialization mismatch, statistical rows and KAT
checks overlap with the cited analyses; see \url{FINDINGS.md}.

\textbf{Recommendation.} The submission does not presently justify
adoption as a general-purpose cryptographic hash. The main concerns are
an applicable subgroup collision-search strategy missing from the
security analysis, complete forgery of several naive keyed
constructions, substantial restrictions on digest use, and unsupported
security arguments. The specification and implementations also disagree
on serialization, and the statistical reports do not support the
statement that every test passed.

The evidence has important limits: this assessment demonstrates
full-parameter forgeries against specified MAC constructions and
recovery of a short preimage, but \textbf{does not exhibit a practical
full-parameter collision, a distinct second preimage, or recovery of a
uniformly random long message from one digest}. Collision-cost
projections below identify their assumptions. The supplied specification
does not assign explicit numerical work factors to each security
property; this report does not infer such claims from digest size alone.

\Needspace{5\baselineskip}

\Needspace{10\baselineskip}

\section{1. Construction and independently checked implementation
behavior}\label{construction-and-independently-checked-implementation-behavior}

\Needspace{6\baselineskip}

Write

\[
A=\begin{pmatrix}-2&1\\4&-3\end{pmatrix},\qquad
B=\begin{pmatrix}-5&2\\-6&2\end{pmatrix},\qquad
M(x_1\ldots x_n)=G_{x_1}\cdots G_{x_n},
\]

where \(G_0=A\), \(G_1=B\), and arithmetic is modulo \(p\). CHAMP-512
uses \(p=2^{128}-15449\); CHAMP-1024 uses \(p=2^{256}-36113\). Both
generators have determinant 2. The digest is \(H(X)=E(M(X))\), where
\(E\) inverts each nonzero entry, leaves zero unchanged, and serializes
the entries in column order. On canonical field encodings, \(E\) is a
publicly invertible bijection. Let \(D=E^{-1}\).

An independent model reproduced all 4,097 short-message vectors for each
variant and agreed with both the reference and optimized C
implementations: 8,194 vectors checked against two implementations
apiece. Eight large all-zero/all-one vectors were independently checked
by matrix exponentiation. The large DRBG-expanded and million-iteration
loop vectors were not rerun in this assessment.
\href{REPRODUCING.md}{Reproduction code and coverage},
\href{evidence/reproduce-full.log}{run transcript}.

The Julia listing on PDF p.~5 writes each field element in big-endian
byte order. All four C implementations and the KATs use little-endian
byte order. Section 5 does not resolve that choice explicitly. This is
an interoperability defect: even the one-bit message \texttt{1} yields
different byte strings. The findings below apply under either
convention; the demonstrations use the C/KAT convention.

\Needspace{5\baselineskip}

\section{2. A subgroup collision-search
strategy}\label{a-subgroup-collision-search-strategy}

Mullan and Tsaban give a subgroup-based approach to short collisions in
matrix homomorphic hashes. Their general-case square-root estimate is
heuristic. It is materially more relevant than inferring an attack
merely from membership in a previously studied family.
\href{https://arxiv.org/abs/1306.5646}{Mullan--Tsaban, §4}.

\Needspace{8\baselineskip}

For CHAMP's implemented pair, the following adaptation works directly
with determinant-2 matrices; no conversion of unequal-length collisions
from an \(SL_2\) construction is needed. Put

\[
P=\begin{pmatrix}1&2\\2&3\end{pmatrix},\quad
P^{-1}AP=\begin{pmatrix}-4&1\\2&-1\end{pmatrix}=C,\quad
P^{-1}BP=\operatorname{diag}(-1,-2)=F.
\]

These are integer identities and hold at both full parameter sets.
Conjugation preserves products and collisions.

\Needspace{7\baselineskip}

First find a positive bitstring \(w\), containing a zero, whose
conjugated product \(T\) is upper triangular. This can be searched by
matching projective points

\[
M(u)^{-1}\infty=M(v)\infty,
\]

where products in this paragraph use \(C,F\), and \(\infty=[1:0]\). A
match makes \(M(uv)\) upper triangular. Only the search computation uses
inverses: the output word is the ordinary positive word \(uv\). There
are \(p+1\) projective points. If the sampled point distributions mix
adequately, two lists of approximately \(\sqrt p\) words suffice.

\Needspace{6\baselineskip}

Now multiply the two upper triangular matrices \(F,T\). For

\[
R=\begin{pmatrix}a&b\\0&d\end{pmatrix},\qquad
S=\begin{pmatrix}e&f\\0&h\end{pmatrix},
\]

\Needspace{6\baselineskip}

direct multiplication gives

\[
RS=SR\quad\Longleftrightarrow\quad f(a-d)=b(e-h).
\]

\Needspace{10\baselineskip}

For nonzero off-diagonal entries, matching the field value \((a-d)/b\)
therefore finds commuting products. Diagonal matrices commute with \(F\)
and can be handled separately. A birthday search on these codes costs
approximately \(\sqrt p\) samples under a corresponding mixing
assumption. Expand the two resulting words back into original CHAMP
bitstrings \(r,s\), and verify \(rs\ne sr\). The resulting collision
satisfies

\[
M(rs)=M(sr),\qquad |rs|=|sr|.
\]

This construction automatically respects the determinant constraint. It
produces an actual collision, not an identity word or a relation
requiring inverse letters in the message.

Our \href{code/subgroup.py}{implementation} produced twelve distinct
equal-length collision pairs at reduced field sizes, using safe primes
congruent to 7 modulo 8. The three runs at each size used the following
total sample counts across the two stages:

\Needspace{9\baselineskip}

{\def\LTcaptype{none} % do not increment counter
\begin{longtable}[]{@{}
  >{\raggedright\arraybackslash}p{(\linewidth - 4\tabcolsep) * \real{0.1700}}
  >{\raggedleft\arraybackslash}p{(\linewidth - 4\tabcolsep) * \real{0.2500}}
  >{\raggedright\arraybackslash}p{(\linewidth - 4\tabcolsep) * \real{0.5800}}@{}}
\toprule\noalign{}
\begin{minipage}[b]{\linewidth}\raggedright
Field bits
\end{minipage} & \begin{minipage}[b]{\linewidth}\raggedleft
Prime
\end{minipage} & \begin{minipage}[b]{\linewidth}\raggedright
Total samples, three seeds
\end{minipage} \\
\midrule\noalign{}
\endhead
\bottomrule\noalign{}
\endlastfoot
16 & 65,063 & 639; 667; 634 \\
20 & 1,048,343 & 1,691; 3,770; 2,357 \\
24 & 16,774,679 & 11,880; 9,224; 8,007 \\
28 & 268,434,263 & 30,771; 33,867; 39,728 \\
\end{longtable}
}

Every pair was recomputed with the original generators and finalization.
A separate verifier recomputed the complete products over the integers
before reducing modulo each prime; it also confirmed that these
particular pairs do not collide at either full parameter set. Complete
messages, digests, seeds and timings are retained in
\href{data/collisions.json}{collision certificates}. These examples test
the mechanism; they do not prove its asymptotic cost or full-size mixing
assumptions.

With words of \(O(\log p)\) symbols at each stage, the demonstrated
approach suggests \(O(\sqrt p\log p)\) field work, approximately
\(\sqrt p\) stored records up to polynomial factors, and messages of
\(O((\log p)^2)\) bits. The leading exponential work factors would be
\textbf{\(2^{64}\) for CHAMP-512 and \(2^{128}\) for CHAMP-1024},
subject to those assumptions. These are neither measured full-parameter
attacks nor established lower bounds. The authors should address this
strategy explicitly before assigning security levels.

\Needspace{5\baselineskip}

\Needspace{10\baselineskip}

\section{3. What follows unconditionally from the
determinant}\label{what-follows-unconditionally-from-the-determinant}

\Needspace{6\baselineskip}

For an \(n\)-bit message,

\[
\det M(X)=2^n\pmod p.
\]

\Needspace{6\baselineskip}

Consequently, outputs at a fixed length occupy at most

\[
S=|SL_2(\mathbb F_p)|=p(p^2-1)
\]

matrices, regardless of which subgroup the generators produce. A proof
that the generators generate a large group would not imply that every
matrix in this coset occurs at a particular length.

\Needspace{16\baselineskip}

{\def\LTcaptype{none} % do not increment counter
\begin{longtable}[]{@{}
  >{\raggedright\arraybackslash}p{(\linewidth - 6\tabcolsep) * \real{0.3000}}
  >{\centering\arraybackslash}p{(\linewidth - 6\tabcolsep) * \real{0.1700}}
  >{\centering\arraybackslash}p{(\linewidth - 6\tabcolsep) * \real{0.1700}}
  >{\raggedright\arraybackslash}p{(\linewidth - 6\tabcolsep) * \real{0.3600}}@{}}
\toprule\noalign{}
\begin{minipage}[b]{\linewidth}\raggedright
Result
\end{minipage} & \begin{minipage}[b]{\linewidth}\centering
CHAMP-512
\end{minipage} & \begin{minipage}[b]{\linewidth}\centering
CHAMP-1024
\end{minipage} & \begin{minipage}[b]{\linewidth}\raggedright
Interpretation
\end{minipage} \\
\midrule\noalign{}
\endhead
\bottomrule\noalign{}
\endlastfoot
Fixed-length range bound & Less than \(2^{384}\) & Less than \(2^{768}\)
& Exact counting bound \\
Classical birthday collision search & \(O(2^{192})\) & \(O(2^{384})\) &
Generic baseline; choose a domain with at least \(2S\) messages \\
Some equal-length collision must exist by & 384 input bits & 768 input
bits & Pigeonhole principle; does not locate the pair \\
Subgroup collision search & About \(2^{64}\) & About \(2^{128}\) &
Leading exponential factors only; mixing assumptions from item 2 \\
\end{longtable}
}

The first three rows are independent of the subgroup heuristic. The
birthday row assumes sampling at a length with at least \(2S\) distinct
possible messages and discarding repeated inputs; it is not a bound for
a nearly injective short-message domain. The pigeonhole lengths are 384
and 768, since the range sizes are already strictly smaller than the
corresponding powers of two.

\Needspace{8\baselineskip}

Let \(q=(p-1)/2\). Proof-enabled primality checks confirmed that both
\(p,q\) are prime, and excluded every larger candidate safe prime at
each size. Since \(p\equiv7\pmod8\), 2 has order \(q\). Thus a collision
implies

\[
|X|\equiv |Y|\pmod q.
\]

This is a valid and useful exclusion of different-length collisions
within a range of width less than \(q\). It does not address same-length
collision search.

Two corrections are needed in the specification's statements on PDF
pp.~4--5. Equality gives \(2^{|X|-|Y|}=1\); allowing \(\pm1\) is an
unnecessarily weak condition, rather than a contradiction. An identity
word has length divisible by \(q\), not necessarily equal to \(q\) or
exactly \(2^{255}\).

For CHAMP-1024, \(A^q=I\) and \(q=2^{255}-18057\). The symbolic messages
\(\epsilon\) and \(0^q\) therefore contradict the literal assertion of
no different-length collisions with difference less than \(2^{255}\).
Matrix exponentiation verifies the identity; the enormous message has
not been expanded, and this is not a practical attack. The correct
threshold is \(q\).

The two variants must also be distinguished. At 256 field bits,
\(\operatorname{ord}(A)=q\) and \(\operatorname{ord}(B)=2q\). At 128
field bits, 17 is a quadratic \textbf{nonresidue}, \(A\) is nonsplit,
and \(A^q\ne I\), while \(\operatorname{ord}(B)=2q\) still holds. The
alternative matrix with trace 5 has the same discriminant 17 and the
same splitting distinction. Arguments requiring split generators
therefore need separate treatment at the two sizes.
\href{evidence/parameters.log}{Parameter proofs and transcript}.

The determinant also exposes length modulo \(q\). If the length lies in
a known interval of width \(L<q\), bounded discrete-log methods recover
it in \(O(\sqrt L)\) group operations, with an appropriate time--memory
tradeoff. Testing a proposed length takes a modular exponentiation.
Unbounded exact length recovery does not follow. Length hiding is an
application requirement, not a universal hash-function requirement.

\Needspace{5\baselineskip}

\Needspace{10\baselineskip}

\section{4. Exact failures of particular keyed
constructions}\label{exact-failures-of-particular-keyed-constructions}

\Needspace{6\baselineskip}

The public combination rule is

\[
H(XY)=E(D(H(X))D(H(Y))).
\]

\Needspace{7\baselineskip}

For a naive prefix MAC \(t=H(Km)\), one known message--tag pair reveals
an equivalent key matrix:

\[
L=D(t)M(m)^{-1}=M(K).
\]

For any chosen replacement \(m'\), the attacker returns \(E(LM(m'))\).
This is universal message forgery after one observed tag, without
recovering the key bits or knowing their length. It works for
replacements of the same length, so requiring a fixed message length
does not fix this construction. A suffix MAC \(H(mK)\) fails similarly
by multiplying by \(M(m)^{-1}\) on the left.

\Needspace{7\baselineskip}

Adding a fixed secret suffix as well as a prefix is insufficient.
Suppose decoded tags have the form \(T(m)=L M(m)R\). Query tags for
\(\epsilon,0,1\). Then

\[
T(0)T(\epsilon)^{-1}=LAL^{-1},\qquad
T(1)T(\epsilon)^{-1}=LBL^{-1}.
\]

Multiplying these known conjugated generators according to a new message
and then multiplying by \(T(\epsilon)\) produces its valid tag. These
attacks were checked against all four shipped C implementations. These
constructions are assessed as potential uses of CHAMP, not as MAC
schemes expressly proposed in the specification. The results are not an
analysis of HMAC or every possible keyed construction.

\Needspace{6\baselineskip}

For public digests, if \(H(X)\) and \(H(Xb)\) are both disclosed, then

\[
D(H(X))^{-1}D(H(Xb))=G_b.
\]

This recovers the appended bit exactly at any prefix length. The
demonstrations recover all 5,000 bits from 5,001 consecutive digests in
the independent model and check recovery of bit 5,000 using each C
implementation. This identifies a confidentiality hazard \textbf{when
such related digests are published}. Streaming computation itself does
not require publishing intermediate digests. This attack therefore does
not refute the distinct single-digest prediction claim in §7.1(2); that
claim still needs a precise input distribution and security experiment.

Public digest combination also does not, by itself, break commitment
binding or file integrity against a protected digest. Changing a
commitment value is not opening the same commitment two ways.
Recomputing a digest for a modified file does not make it equal an
authenticated original digest. Any claimed protocol failure must specify
what the adversary may change and what the verifier trusts.

\Needspace{5\baselineskip}

\Needspace{11\baselineskip}

\section{5. Short-input recovery and preimage-search
bounds}\label{short-input-recovery-and-preimage-search-bounds}

\Needspace{7\baselineskip}

For a target matrix \(T\) known to hash an \(m\)-bit input, split a
candidate as \(uv\), with lengths as close as possible to \(m/2\). Store
\(M(u)\), and look up

\[
T M(v)^{-1}.
\]

This finds an \(m\)-bit preimage using \(O(2^{m/2})\) stored values and
operations up to polynomial factors. It works directly modulo \(p\); no
lifting to integers or assumption about entry growth is required. Our
independent implementation recovered a 40-bit input from its CHAMP-1024
digest in about 6 seconds on the recorded environment. That timing
describes this run, not a universal practical cutoff. A second
demonstration succeeds after integer wrap at a reduced prime.
\href{evidence/reproduce-full.log}{Transcript}.

The distinction from exhaustive guessing is significant: a generic
random hash of a uniformly chosen 40-bit input normally requires
searching roughly \(2^{40}\) possibilities. Nevertheless, this example
does not recover arbitrary 128-bit or long high-entropy inputs in
practical time.

Integer exposure is a separate observation. Recovering centered matrix
entries exposes the exact integer product only when all original entries
lie in \([-\lfloor p/2\rfloor,\lfloor p/2\rfloor]\). Random samples
cannot establish a universal no-wrap threshold. In particular, all-zero
inputs leave this interval at \textbf{59 bits for CHAMP-512 and 117 bits
for CHAMP-1024}, providing a worst-case counterexample to any larger
universal centered-lifting threshold. Exposure of a matrix is also not
automatically efficient recovery of its word factorization.

Two further distinctions are essential:

\begin{itemize}
\tightlist
\item
  A preimage algorithm may return the original message. It is not a
  second-preimage attack unless it finds a different message. All 4,096
  twelve-bit messages have distinct products at both full parameter sets
  in our exhaustive check; the meet-in-the-middle search in that domain
  returns the original. The determinant also excludes different feasible
  lengths for a second preimage of these targets.
\item
  A saturated \(O(\sqrt S)\) target-preimage estimate needs sufficiently
  long halves and suitably spread forward/backward distributions, so
  that sampled lists intersect at the target. Range cardinality alone
  does not prove this estimate for every length or target. A toy success
  does not supply the missing hypothesis.
\end{itemize}

Quantum estimates need a specified resource model. In particular, the
two split searches form a claw-finding problem on domains of size about
\(2^{m/2}\). Applying Tani's algorithm yields \(O(2^{m/3})\) quantum
queries to these product functions, with reversible evaluation and
substantial quantum storage requirements. This is an inference from that
algorithm, not a tested quantum implementation.
\href{https://arxiv.org/abs/0708.2584}{Tani, Corollary 9}. Standard BHT
collision estimates also require their oracle and balance assumptions to
be stated.
\href{https://arxiv.org/abs/quant-ph/9705002}{Brassard--Høyer--Tapp}. No
optimal quantum security level is established here.

\Needspace{5\baselineskip}

\Needspace{12\baselineskip}

\section{6. Statistical appearance, security arguments, and graph
terminology}\label{statistical-appearance-security-arguments-and-graph-terminology}

\Needspace{8\baselineskip}

A single digest at a known length admits a strong structural test:
reject noncanonical field words, decode, and test determinant \(2^n\).
CHAMP always passes. A uniformly random \(4b\)-bit string passes with
probability exactly

\[
\frac{p(p^2-1)}{2^{4b}}\approx 2^{-b},\qquad b=\operatorname{bitlength}(p).
\]

Thus statistical test-suite success would not make these digests behave
as uniform random outputs. A three-query concatenation test is also
possible; conditional on two valid decoded outputs, matching their
encoded product in a fresh random \(k\)-bit oracle response has
probability \(2^{-k}\). These properties prevent directly transferring a
random-oracle argument without examining its use of the hash. They do
not imply that every protocol with a random-oracle proof is insecure.

The supplied NIST reports themselves contain four marked failures:

\Needspace{9\baselineskip}

{\def\LTcaptype{none} % do not increment counter
\begin{longtable}[]{@{}
  >{\raggedright\arraybackslash}p{(\linewidth - 4\tabcolsep) * \real{0.4000}}
  >{\raggedleft\arraybackslash}p{(\linewidth - 4\tabcolsep) * \real{0.1000}}
  >{\raggedright\arraybackslash}p{(\linewidth - 4\tabcolsep) * \real{0.5000}}@{}}
\toprule\noalign{}
\begin{minipage}[b]{\linewidth}\raggedright
Report
\end{minipage} & \begin{minipage}[b]{\linewidth}\raggedleft
Line
\end{minipage} & \begin{minipage}[b]{\linewidth}\raggedright
Marked result
\end{minipage} \\
\midrule\noalign{}
\endhead
\bottomrule\noalign{}
\endlastfoot
Primary pair, 512-bit digest & 156 & NonOverlappingTemplate, 95/100 \\
Primary pair, 1024-bit digest & 63 & NonOverlappingTemplate, 94/100 \\
Primary pair, 1024-bit digest & 123 & NonOverlappingTemplate, 95/100 \\
Alternative pair, 512-bit digest & 186 & RandomExcursionsVariant,
56/60 \\
\end{longtable}
}

Exact filenames, source hashes and the four flagged rows are recorded in
the \href{PROVENANCE.md}{source inventory} and
\href{data/nist-flagged-rows.json}{row extracts}. The two primary-pair
reports state a minimum acceptable ordinary-test count of approximately
96/100. The authors should reconcile these records with the claim that
all tests passed. Some failures are expected when many statistical tests
are applied; their presence alone does not prove an exploitable bias.
More fundamentally, statistical tests cannot replace cryptanalysis.
\href{https://csrc.nist.gov/pubs/sp/800/22/r1/upd1/final}{NIST SP 800-22
Rev.~1a}.

Testing 100 streams of \(10^6\) bits uses 12.5 MB. This is compatible
with generating at least 100 MB and testing a subset. The exact
generated amount, tested subset, counter encoding and testing procedure
should be documented.

The length-attack defense on PDF p.~7 does not establish resistance. A
large canonical representative of \(2^{-1}\) says little about whether
applying the \textbf{correct} inverse factor recovers a smaller integer
predecessor. Nor does the sign of an average-matrix eigenvalue control
the norm of an individual product. Replacing both generators by their
negatives changes an \(n\)-factor product only by \((-1)^n\), leaving
every absolute-entry norm unchanged, while reversing the signs of the
average matrix's eigenvalues. Furthermore,
\(\mathbb E[M_n]=((A+B)/2)^n\) concerns signed entries, not
\(\mathbb E[\|M_n\|]\), typical growth, or a predictor's success
probability. The displayed plots do not supply those missing analyses.

The term \emph{girth} must be defined. A positive collision gives two
directed paths with a common endpoint, but closing one path by the
reverse of the other uses inverse edges. It bounds a cycle in the
underlying undirected graph, not necessarily a directed identity cycle.
Indeed, for the CHAMP-1024 primary pair the directed girth is exactly
\(q\): determinants rule out shorter positive identity words, and
\(A^q=I\) supplies one. This can coexist with same-length collisions
forced by counting at 768 bits. Neither a large directed girth nor a
counting upper bound for a different graph resolves the actual
collision-search problem.

The integer-lifting discussion on PDF p.~6 also needs care with signed
entries. Two different integers can both have absolute value less than
\(p\) and still be equal modulo \(p\): 1 and \(1-p\) are an example. A
sufficient no-aliasing condition is that each coordinate lies in a
common interval of width less than \(p\), such as the centered interval
used above. Freeness over the integers and a valid coordinate bound are
separate requirements for such a modular collision exclusion.

A bounded search without an integer relation does not prove freeness.
Conversely, lack of a freeness proof is not itself a collision attack. A
credible security analysis must address the actual parameters and
message domain.

\Needspace{5\baselineskip}

\section{7. Implementation and feature
claims}\label{implementation-and-feature-claims}

All four \texttt{CryptHash} implementations ignore
\texttt{digest\_len\_bits}. For example, asking the CHAMP-1024
implementation for 512 bits still returns success and writes 1,024 bits.
A guarded-buffer test confirms this behavior without making an
out-of-bounds access. The API should reject unsupported lengths and
document output capacity.

All four implementations lazily initialize global byte tables using an
unsynchronized ready flag. Concurrent first calls can race on these
shared objects under the C memory model. This is a source-level finding,
not a reproduced erroneous digest. Constant tables or properly
synchronized initialization would remove that concern.

The implementation is not established to be constant-time. Its memory
accesses use message bytes as table indices, and arithmetic contains
data-dependent control flow. A public exponent fixes an exponentiation
schedule, not the timing of every arithmetic operation. Hash inputs may
be confidential in some applications. No timing or cache attack is
demonstrated here. Implementations intended for confidential inputs need
a separate side-channel assessment.

Streaming without knowing the final length is already supported by
SHA-family hashing: complete blocks can be processed as they arrive and
padding added at the end. CHAMP's public digest-combination rule is a
different functionality. The feature comparison in §§2 and 4 should
distinguish them.
\href{https://nvlpubs.nist.gov/nistpubs/FIPS/NIST.FIPS.180-4.pdf}{FIPS
180-4, §5.1}.

The paper's own Table 1 reports lower single-thread throughput than
SHA-512; exploiting more cores is a separate resource comparison. No new
throughput benchmark is claimed here. Breaking public homomorphism is
not a necessary condition for collision resistance, and conventional
streaming or parallel hashing does not inherently require publishing a
composable final digest.

\textbf{Requested revisions before reconsideration.} State separate
collision, preimage and second-preimage targets; specify input
distributions, message limits and classical/quantum resource models;
address the subgroup strategy and its mixing assumptions; formalize the
single-digest bit-prediction claim; limit proposed uses to analyzed
constructions; correct the determinant statements and serialization;
reconcile the statistical reports; and repair or document the API and
concurrency behavior. Any revised finalization must be analyzed as part
of the construction: simply hashing the matrix again cannot remove
collisions already present in the matrix product.

The accompanying \href{REPRODUCING.md}{reproduction guide} maps the
computations to source code, complete collision certificates and
recorded transcripts. The \href{PROVENANCE.md}{source inventory}
identifies the exact reviewed inputs and version boundaries. The
\href{FINDINGS.md}{findings index} separates algebraic results,
experiments, projections and source-level observations.

\Needspace{8\baselineskip}

\textbf{AI-use disclosure.} Moonshot AI's kimi-k3 was used for the main
review. OpenAI's GPT-6 Astra was used for polishing and preparing the
publication package, including additional technical analysis,
cross-checking of claims, and development of reproducibility code.
Mounir IDRASSI is responsible for the final report and its conclusions.

\textbf{License.} Copyright © 2026 Mounir IDRASSI for the original
contributions. This report and its original results are licensed under
\href{https://creativecommons.org/licenses/by/4.0/}{Creative Commons
Attribution 4.0 International (CC BY 4.0)}. The accompanying code is
licensed under the \href{LICENSE}{MIT License}. Third-party material
retains its existing terms; see \url{LICENSING.md} for the scope.

\Needspace{5\baselineskip}

\section{8. References and
reproducibility}\label{references-and-reproducibility}

\begin{enumerate}
\def\labelenumi{\arabic{enumi}.}
\item
  CHAMP submitters. \emph{CHAMP: Cayley HAshing with Matrix Products ---
  Algorithm Specifications}. Supplied ten-page specification, undated.
  Reviewed file: \texttt{CHAMP\ Algorithm\ Specifications.pdf}. SHA-256,
  split over two lines for readability:

\begin{Shaded}
\begin{Highlighting}[]
\NormalTok{c1201fe48fe45033c9a97e39ff83928e}
\NormalTok{683ac1d69703ec31987752f6f7dfaa69}
\end{Highlighting}
\end{Shaded}
\item
  Alexander Demin, Alexey Ovchinnikov and Vladimir Shpilrain.
  \emph{CHAMP: Cayley HAshing with Matrix Products}. Distinct
  twelve-page
  \href{https://shpilrain.ccny.cuny.edu/Cayley\%20hash\%20paper.pdf}{author-hosted
  preprint}, retrieved 21 September 2026. Version identification only;
  see the scope note above.
\item
  Ciaran Mullan and Boaz Tsaban. \emph{SL2 homomorphic hash functions:
  Worst case to average case reduction and short collision search}.
  \href{https://arxiv.org/abs/1306.5646v3}{arXiv:1306.5646v3}, 23
  December 2015, especially §4. Journal DOI:
  \href{https://doi.org/10.1007/s10623-015-0129-8}{\nolinkurl{10.1007/s10623-015-0129-8}}.
\item
  Seiichiro Tani. \emph{Claw Finding Algorithms Using Quantum Walk}.
  \href{https://arxiv.org/abs/0708.2584v2}{arXiv:0708.2584}, version 2,
  3 March 2008, Corollary 9. \emph{Theoretical Computer Science}
  410(50), 5285--5297 (2009). DOI:
  \href{https://doi.org/10.1016/j.tcs.2009.08.030}{\nolinkurl{10.1016/j.tcs.2009.08.030}}.
\item
  Gilles Brassard, Peter Høyer and Alain Tapp. \emph{Quantum Algorithm
  for the Collision Problem}.
  \href{https://arxiv.org/abs/quant-ph/9705002}{arXiv:quant-ph/9705002},
  1997. \emph{LATIN 1998}, LNCS 1380, 163--169. DOI:
  \href{https://doi.org/10.1007/BFb0054319}{\nolinkurl{10.1007/BFb0054319}}.
\item
  NIST. \emph{A Statistical Test Suite for Random and Pseudorandom
  Number Generators for Cryptographic Applications}.
  \href{https://csrc.nist.gov/pubs/sp/800/22/r1/upd1/final}{SP 800-22
  Rev.~1a}, 2010. DOI:
  \href{https://doi.org/10.6028/NIST.SP.800-22r1a}{\nolinkurl{10.6028/NIST.SP.800-22r1a}}.
\item
  NIST. \emph{Secure Hash Standard}.
  \href{https://nvlpubs.nist.gov/nistpubs/FIPS/NIST.FIPS.180-4.pdf}{FIPS
  PUB 180-4}, August 2015, §5.1. DOI:
  \href{https://doi.org/10.6028/NIST.FIPS.180-4}{\nolinkurl{10.6028/NIST.FIPS.180-4}}.
\item
  Mounir IDRASSI. \emph{CHAMP: technical assessment of the submitted
  specification}, accompanying reproducibility package, version 1.0.1,
  22 September 2026. The archive contains the independent model,
  reduced-parameter search, separately implemented certificate verifier,
  parameter proofs and transcripts. SHA-256 file checksums accompany the
  release; they establish file integrity, not authorship or
  cryptanalytic correctness.
\item
  Markku-Juhani O. Saarinen.
  \emph{\href{https://ngcc.dev/reports/hash-04.html}{CHAMP (hash-04)}}.
  ngcc.dev, 21 September 2026; retrieved 22 September 2026.
\item
  cuihr26 (signed ``Tsinghua Hash Lab'').
  \emph{\href{https://list.niccs.org.cn/archives/list/crypthashforum@list.niccs.org.cn/message/QNMIAAKIPHG5FEK7QZEE2ZYBOWQHCZUE/}{Re:
  Round1: Public Comment: CHAMP}}. CryptHash Forum, 22 September 2026;
  retrieved the same day.
\item
  Cryptanalysts001 (signed ``ISCAS'').
  \emph{\href{https://list.niccs.org.cn/archives/list/crypthashforum@list.niccs.org.cn/message/6LSOFTTMTDWFWWRAZDREDYJNG5NXDVEC/}{Re:
  Round1: Public Comment: CHAMP}}. CryptHash Forum, 22 September 2026;
  retrieved the same day.
\end{enumerate}

\end{document}
